\section{TarMAC: Targeted Multi-Agent Communication}
\label{sec:approach}

\begin{figure*}[t]
    \centering
    \includegraphics[width=\textwidth]{figures/main/approach_v5.png}
    \vspace{2pt}
    \caption{Overview of our multi-agent architecture with targeted communication.
    Left: At every timestep, each agent policy gets a local observation $\omega_i^t$
    and aggregated message $c_i^t$ as input, and predicts an environment action $a_i^t$
    and a targeted communication message $m_i^t$. Right: Targeted communication
    between agents is implemented as a signature-based soft attention mechanism.
    Each agent broadcasts a message $m_i^t$ consisting of a signature $k_i^t$,
    which can be used to encode agent-specific information and
    a value $v_i^t$, which contains the actual message. At the next timestep, each receiving agent
    gets as input a convex combination of message values, where the
    attention weights are obtained by a dot product between sender's signature $k_i^t$
    and a query vector $q_j^{t+1}$ predicted from the receiver's hidden state.}
    \vspace{-15pt}
\end{figure*}

We now describe our multi-agent communication architecture in detail.
Recall that we have $N$ agents with policies $\{\pi_1, ..., \pi_N\}$, respectively
parameterized by $\{\theta_1, ..., \theta_N\}$, jointly performing a cooperative task.
At every timestep $t$, the $i$th agent for all $i \in \{1, ..., N\}$ sees a local
observation $\omega_i^t$, and must select a discrete environment action $a_i^t \sim \pi_{\theta_i}$
and send a continuous communication message $m_i^t$, received by other agents at the next timestep,
in order to maximize global reward $r_t \sim R$. Since no agent has access to the
underlying complete state of the environment $s_t$, there is incentive in communicating
with each other and being mutually helpful to do better as a team.

\textbf{Policies and Decentralized Execution.} Each agent is essentially modeled as a
Dec-POMDP augmented with communication.  Each agent's policy $\pi_{\theta_i}$
is implemented as a $1$-layer Gated Recurrent Unit~\citep{cho_emnlp14}.
At every timestep, the local observation $\omega_i^t$ and a vector $c_i^t$
aggregating messages sent by all agents at the previous timestep (described in
more detail below) are used to update the hidden state $h_i^t$ of the GRU,
which encodes the entire message-action-observation history up to time $t$.
From this internal state representation, the agent's policy
$\pi_{\theta_i} \left(a_i^t \, | \, h_i^t \right)$ predicts a categorical
distribution over the space of actions, and another output head produces an
outgoing message vector $m_i^t$. Note that for our experiments,
agents are symmetric and policy parameters are shared across agents,
\ie $\theta_1 = ... = \theta_N$.
This considerably speeds up learning.

\textbf{Centralized Critic.} Following prior work~\citep{lowe_nips17,foerster_aaai18},
we operate under the centralized learning and decentralized execution
paradigm wherein during training, a centralized Critic
guides the optimization of individual agent policies. The Critic takes as
input predicted actions $\{ a_1^t, ..., a_N^t \}$ and internal state
representations $\{h_1^t, ..., h_N^t \}$ from all agents to estimate the
joint action-value $\hat Q_t$ at every timestep. The centralized Critic
is learned by temporal difference~\citep{sutton_book98} and the gradient
of the expected return $J(\theta_i) = \mathbb{E}[R]$ with respect to
policy parameters is approximated by:

\resizebox{1. \columnwidth}{!}
{
$
    \nabla_{\theta_{i}} J(\theta_i) =
        \mathbb{E}
            \left[ \nabla_{\theta_{i}} \log \pi_{\theta_{i}} (a_i^t | h_i^t)
                \,\, \hat{Q}_{t} (h_1^t, ..., h_N^t, a_t^1, ..., a_t^N) \right].
$
}


Note that compared to an individual Critic $\hat{Q}_i(h_i^t, a_i^t)$ per agent,
having a centralized Critic leads to considerably lower variance
in policy gradient estimates since it takes into account actions from all agents.
At test time, the Critic is not needed and policy execution is
fully decentralized.

\textbf{Targeted, Multi-Round Communication.}
Establishing complex collaboration strategies requires targeted communication
\ie the ability to address specific messages to specific agents, as well as
multi-round communication \ie multiple rounds of back-and-forth interactions between
agents. We use a signature-based soft-attention mechanism in our communication
structure to enable targeting. Each message $m_i^t$ consists of $2$ parts:
a signature $k_i^t \in \mathbb{R}^{d_k}$ to encode properties of intended recipients
and a value $v_i^t \in \mathbb{R}^{d_v}$:
\begin{equation}\label{eq:message}
m_i^t = [\,\,\,\overbracket{\,\,\,\,\,\, k_i^t \,\,\,\,\,\,}^{\textnormal{signature}} \quad
    \underbracket{\,\,\, v_i^t \,\,\,}_{\textnormal{value}}\,\,\,] \;.
\end{equation}
At the receiving end, each agent (indexed by $j$) predicts a query
vector $q_j^{t+1} \in \mathbb{R}^{d_k}$ from its hidden state $h_j^{t+1}$,
which is used to compute a dot product with signatures of all $N$ messages.
This is scaled by $1 / \sqrt{d_k}$ followed by a softmax to
obtain attention weights $\alpha_{ji}$ for each incoming message:
\begin{equation}\label{eq:alpha}
\mathbf{\alpha}_{j} = \textnormal{softmax}\left[
\frac{{q_j^{t+1}}^T {k_1^t}}{\sqrt{d_k}} \,\, ... \,\,
    \underbracket{\frac{ {{q_j^{t+1}}^T {k_i^t}}}{\sqrt{d_k}}}_{\alpha_{ji}} \,\, ... \,\, \frac{{q_j^{t+1}}^T {k_N^t}}{\sqrt{d_k}} \right] \quad \quad \quad
\end{equation}

\begin{figure*}[t!]
    \centering
    \includegraphics[width=\textwidth]{figures/main/shapes_example_3x3_rrgb.jpg}
    \vspace{-15pt}
    \caption{Visualizations of learned targeted communication in SHAPES. Figure best viewed in color.
    $4$ agents have to find $[$\texttt{red}, \texttt{red}, \texttt{green}, \texttt{blue}$]$ respectively.
    $t=1$: inital spawn locations;
    $t=2$: $4$ was on \texttt{red} at $t=1$ so $1$ and $2$ attend to messages from $4$ since they have to find \texttt{red}. $3$ has found its goal (\texttt{green}) and is self-attending;
    $t=6$: $4$ attends to messages from $2$ as $2$ is on $4$'s target -- \texttt{blue};
    $t=8$: $1$ finds \texttt{red}, so $1$ and $2$ shift attention to $1$;
    $t=21$: all agents are at their respective goal locations and primarily self-attending.}
    \label{fig:shapes_3x3}
    \vspace{-15pt}
\end{figure*}
used to compute
$c_j^{t+1}$, the input message for agent $j$ at $t+1$:
\begin{equation}\label{eq:attended}
c_j^{t+1} = \sum_{i=1}^N \alpha_{ji} v_i^t.
\end{equation}
Intuitively, attention weights are high when both sender and receiver
predict similar signature and query vectors respectively.
Note that \Eqref{eq:alpha} also includes $\alpha_{ii}$ corresponding
to the ability to \emph{self-attend}~\citep{vaswani_nips17}, which we empirically found
to improve performance, especially in situations when an agent has found
the goal in a coordinated navigation task and all it is required to
do is stay at the goal, so others benefit from attending to this agent's
message but return communication is not necessary.
Note that the targeting mechanism in our formulation is implicit
\ie agents implicitly encode properties without addressing specific recipients.
For example, in a self-driving car network, a particular message may be for
``cars travelling on the west to east road" (implicitly encoding properties)
as opposed to specifically for ``car 2'' (explicit addressing).

For multi-round communication, aggregated message vector $c_j^{t+1}$ and
internal state $h_j^t$ are first used to predict the next internal state
${h'}_j^t$ taking into account the first round:
\begin{equation}\label{eq:nextstate}
{h'}_j^t = \tanh\left( W_{h\rightarrow h'} [ \,\, c_j^{t+1} \,\, \Vert \,\, h_j^t \,\, ]  \right).
\end{equation}
Next, the updated hidden state ${h'}_j^t$ is used to predict signature, query, value
followed by repeating Equations~\refeq{eq:message}-\refeq{eq:nextstate} for
multiple rounds until we get a final aggregated message vector $c_j^{t+1}$
to be used as input at the next timestep.
Number of rounds of communication is treated as a hyperparameter.

Our entire communication architecture is differentiable, and
message vectors are learnt through backpropagation.

